% SDU Thesis 配置模板
% 本文件提供了常用的配置选项模板，用户可以根据需要复制和修改

% ==================== 个人信息配置 ====================
% 复制以下内容到您的配置文件中，并修改相应信息

% 论文基本信息
\Title{您的论文标题}
\Author{您的姓名}
\StudentID{您的学号}
\Department{您的学院}
\Major{您的专业}
\Grade{您的年级}
\Teacher{指导教师姓名}

% ==================== 颜色主题配置 ====================
% 选择一个颜色主题，取消注释相应部分

% 蓝色主题
% \definecolor{themecolor}{RGB}{0,102,204}
% \definecolor{accentcolor}{RGB}{0,153,255}

% 绿色主题
% \definecolor{themecolor}{RGB}{0,153,76}
% \definecolor{accentcolor}{RGB}{34,139,34}

% 红色主题
% \definecolor{themecolor}{RGB}{220,20,60}
% \definecolor{accentcolor}{RGB}{255,69,0}

% 紫色主题
% \definecolor{themecolor}{RGB}{138,43,226}
% \definecolor{accentcolor}{RGB}{186,85,211}

% ==================== 字体配置 ====================
% 根据您的系统和需求选择字体配置

% Windows 系统默认字体（推荐）
% \setmainfont{Times New Roman}
% \newCJKfontfamily\songti[Path=assets/fonts/]{SimSun.ttc}
% \newCJKfontfamily\heiti[Path=assets/fonts/]{SimHei.ttf}
% \newCJKfontfamily\kaiti[Path=assets/fonts/]{SimKai.ttf}

% macOS 系统字体
% \setmainfont{Times New Roman}
% \setCJKmainfont{Songti SC}
% \setCJKsansfont{Heiti SC}
% \setCJKmonofont{Kaiti SC}

% Linux 系统字体
% \setmainfont{Liberation Serif}
% \setCJKmainfont{Noto Serif CJK SC}
% \setCJKsansfont{Noto Sans CJK SC}

% ==================== 页面布局配置 ====================
% 选择适合的页面布局

% 标准布局（推荐）
% \geometry{
%     a4paper,
%     left=3cm,
%     right=3cm,
%     top=2.5cm,
%     bottom=2.5cm,
% }

% 紧凑布局
% \geometry{
%     a4paper,
%     left=2.5cm,
%     right=2.5cm,
%     top=2cm,
%     bottom=2cm,
% }

% 宽松布局
% \geometry{
%     a4paper,
%     left=3.5cm,
%     right=3.5cm,
%     top=3cm,
%     bottom=3cm,
% }

% ==================== 章节样式配置 ====================
% 选择章节标题样式

% 传统学术风格（默认）
% \ctexset{
%     chapter={
%       format=\centering\zihao{3}\allbfhei,
%       name={第,章},
%       beforeskip=21.6pt,
%       afterskip=18pt,
%       fixskip=true,
%      },
%     section={
%       format= \zihao{4}\allbfhei,
%       beforeskip=18pt,
%       afterskip=18pt,
%       fixskip=true,
%      }
% }

% 现代简约风格
% \ctexset{
%     chapter={
%       format=\raggedright\zihao{2}\bfseries\color{themecolor},
%       name={},
%       number=\arabic{chapter},
%       beforeskip=0pt,
%       afterskip=30pt,
%       fixskip=true,
%      },
%     section={
%       format=\raggedright\zihao{3}\bfseries\color{themecolor},
%       beforeskip=24pt,
%       afterskip=12pt,
%       fixskip=true,
%      }
% }

% ==================== 图表标题配置 ====================
% 自定义图表标题样式

% 简洁风格
% \captionsetup[figure]{
%     font=small,
%     labelfont=bf,
%     textfont=it,
%     labelsep=period,
%     justification=centering,
% }
% 
% \captionsetup[table]{
%     font=small,
%     labelfont=bf,
%     textfont=normal,
%     labelsep=period,
%     justification=centering,
% }

% 学术风格
% \captionsetup[figure]{
%     font={small,stretch=1.2},
%     labelfont={bf,color=themecolor},
%     textfont=normal,
%     labelsep=colon,
%     justification=justified,
% }

% ==================== 超链接配置 ====================
% 自定义超链接样式

% 彩色链接（屏幕阅读）
% \hypersetup{
%     colorlinks=true,
%     linkcolor=themecolor,
%     citecolor=accentcolor,
%     urlcolor=themecolor,
% }

% 黑白链接（打印版本）
% \hypersetup{
%     colorlinks=false,
%     linkbordercolor=white,
%     citebordercolor=white,
%     urlbordercolor=white,
% }

% ==================== 代码样式配置 ====================
% 自定义代码高亮样式

% Python 代码样式
% \lstdefinestyle{python}{
%     language=Python,
%     backgroundcolor=\color{gray!5},
%     commentstyle=\color{green!60!black},
%     keywordstyle=\color{blue},
%     numberstyle=\tiny\color{gray},
%     stringstyle=\color{orange},
%     basicstyle=\ttfamily\small,
%     breaklines=true,
%     numbers=left,
%     frame=single,
%     rulecolor=\color{gray!30},
% }

% C++ 代码样式
% \lstdefinestyle{cpp}{
%     language=C++,
%     backgroundcolor=\color{blue!5},
%     commentstyle=\color{green!50!black},
%     keywordstyle=\color{blue!80!black},
%     numberstyle=\tiny\color{gray},
%     stringstyle=\color{red!60!black},
%     basicstyle=\ttfamily\footnotesize,
%     breaklines=true,
%     numbers=left,
%     frame=leftline,
%     framerule=2pt,
%     rulecolor=\color{blue!30},
% }

% ==================== 自定义命令 ====================
% 常用的自定义命令

% 强调命令
% \newcommand{\important}[1]{\textcolor{themecolor}{\textbf{#1}}}
% \newcommand{\highlight}[1]{\colorbox{yellow!30}{#1}}
% \newcommand{\note}[1]{\textcolor{gray}{\textit{#1}}}

% 引用命令
% \newcommand{\figref}[1]{图~\ref{#1}}
% \newcommand{\tabref}[1]{表~\ref{#1}}
% \newcommand{\eqnref}[1]{式~(\ref{#1})}
% \newcommand{\chapref}[1]{第~\ref{#1}~章}
% \newcommand{\secref}[1]{第~\ref{#1}~节}

% 数学命令
% \newcommand{\vect}[1]{\boldsymbol{#1}}
% \newcommand{\mat}[1]{\mathbf{#1}}
% \newcommand{\set}[1]{\mathcal{#1}}

% ==================== 自定义环境 ====================
% 常用的自定义环境

% 定理环境
% \newtheoremstyle{mythm}
%     {12pt}{12pt}{\itshape}{0pt}{\bfseries\color{themecolor}}{.}{.5em}{}
% \theoremstyle{mythm}
% \newtheorem{theorem}{定理}[chapter]
% \newtheorem{lemma}[theorem]{引理}
% \newtheorem{definition}[theorem]{定义}

% 示例环境
% \newenvironment{example}[1][]{
%     \par\medskip\noindent
%     \textbf{\color{themecolor}示例%
%     \if\relax\detokenize{#1}\relax\else\ (#1)\fi:}
%     \rmfamily
% }{
%     \par\medskip
% }

% 注意事项环境
% \RequirePackage{tcolorbox}
% \newenvironment{notice}[1][注意]{
%     \begin{tcolorbox}[
%         colback=yellow!10,
%         colframe=orange!80,
%         title=#1,
%         fonttitle=\bfseries,
%     ]
% }{
%     \end{tcolorbox}
% }

% ==================== 使用说明 ====================
% 1. 复制需要的配置到您的自定义配置文件中
% 2. 取消注释（删除行首的 %）
% 3. 根据需要修改参数值
% 4. 在主文档中引入您的配置文件：
%    \input{config/custom/your-config.tex}

% ==================== 配置验证 ====================
% 编译测试清单：
% □ 基础编译无错误
% □ 字体显示正常
% □ 颜色主题正确
% □ 章节样式符合预期
% □ 图表标题格式正确
% □ 超链接工作正常
% □ 代码高亮正确
% □ 自定义命令有效